\documentclass[a4paper,12pt]{article}

\usepackage[T2A]{fontenc}
\usepackage[utf8]{inputenc}
\usepackage[russian]{babel}

\usepackage{amsmath,amssymb,amsthm,mathtools}
\usepackage{helvet}
\renewcommand{\familydefault}{\sfdefault}
\usepackage{icomma}
\usepackage{geometry}
\usepackage{microtype}
\usepackage{tikz}
\usepackage{tkz-euclide}
\usepackage{pgfplots}
\pgfplotsset{compat=1.18}
\usepackage{tabularx,booktabs,longtable}
\usepackage{enumitem}
\usepackage{hyperref}
\usepackage{parskip}
\usepackage{fancyhdr}
\usepackage{setspace}
\usepackage{xcolor}

\geometry{
  a4paper,
  left=18mm,
  right=18mm,
  top=18mm,
  bottom=20mm
}

\definecolor{accent}{HTML}{008080}
\definecolor{darkaccent}{HTML}{004D4D}
\definecolor{textgray}{HTML}{333333}
\definecolor{softgray}{HTML}{6E7777}
\definecolor{lightaccent}{HTML}{DCECEC}

\hypersetup{
  colorlinks=true,
  linkcolor=darkaccent,
  urlcolor=accent,
  pdfauthor={Лёвин Артём Александрович},
  pdftitle={Смешанные уравнения и неравенства}
}

\onehalfspacing
\setlength{\parindent}{0pt}
\setlength{\parskip}{6pt}
\setlist[itemize]{leftmargin=7mm,itemsep=3pt,topsep=4pt}
\setlist[enumerate]{leftmargin=8mm,itemsep=5pt,topsep=4pt}
\renewcommand{\arraystretch}{1.35}

\pagestyle{fancy}
\fancyhf{}
\lhead{\textcolor{textgray}{Уравнения и неравенства}}
\rhead{\textcolor{textgray}{Николь Саркисьянц}}
\cfoot{\textcolor{softgray}{\thepage}}
\renewcommand{\headrulewidth}{0.4pt}
\renewcommand{\headrule}{%
  \hbox to\headwidth{%
    \color{lightaccent}\leaders\hrule height \headrulewidth\hfill}}

\newcommand{\checkitem}{\item[$\square$]}
\newcommand{\important}[1]{\textcolor{darkaccent}{\textbf{#1}}}
\newcommand{\R}{\mathbb{R}}
\newcommand{\Z}{\mathbb{Z}}

\newcommand{\sectionrule}{%
  \vspace{-2mm}
  {\color{lightaccent}\hrule height 1pt}
  \vspace{3mm}
}

\begin{document}

%----------------------------------------------------------------------
% Титульный лист
%----------------------------------------------------------------------

\begin{titlepage}
\thispagestyle{empty}

\begin{center}

\vspace*{25mm}

{\color{accent}\rule{34mm}{1.2pt}}

\vspace{12mm}

{\Large\textcolor{softgray}{ЧЕК-ЛИСТ}}

\vspace{8mm}

{\Huge\bfseries\textcolor{darkaccent}{Смешанные уравнения\\[3mm]
и неравенства}}

\vspace{7mm}

{\Large
ОДЗ, рационализация, метод интервалов\\
и замена переменной}

\vspace{25mm}

{\large
Лёвин Артём Александрович\\[2mm]
эксклюзивно для Николь Саркисьянц}

\vfill

{\large\today}

\vspace{18mm}

\begin{tikzpicture}[x=1cm,y=1cm]
  \draw[color=lightaccent,line width=5pt]
    (-6,0)
    .. controls (-4.1,1.2) and (-2.3,-1.1) .. (0,0)
    .. controls (2.2,1.1) and (4.0,-1.0) .. (6,0);
  \draw[color=accent,line width=1.2pt]
    (-6,0)
    .. controls (-4.1,1.2) and (-2.3,-1.1) .. (0,0)
    .. controls (2.2,1.1) and (4.0,-1.0) .. (6,0);
\end{tikzpicture}

\vspace{6mm}

\end{center}
\end{titlepage}

%----------------------------------------------------------------------
% Основной чек-лист
%----------------------------------------------------------------------

\section*{\textcolor{darkaccent}{1. Главная стратегия}}
\sectionrule

Смешанная задача может одновременно содержать логарифмы, степени,
тригонометрические функции и многочлены. Решение строится
последовательно: сначала ограничения, затем упрощение структуры,
после этого стандартный метод.

\begin{enumerate}
  \item \important{Выпишите ОДЗ} всех логарифмов, знаменателей и корней.
  \item \important{Определите внешнюю структуру}: логарифмическую,
        показательную, рациональную или тригонометрическую.
  \item \important{Упростите выражение} с помощью формул.
  \item \important{Сведите задачу} к знакомому уравнению, неравенству
        или системе.
  \item \important{Решите полученную задачу} методом интервалов,
        заменой переменной или отбором на окружности.
  \item \important{Сверьте результат с ОДЗ}.
\end{enumerate}

\begin{itemize}
  \checkitem Я сначала ищу ограничения, если вижу логарифм.
  \checkitem Я не сокращаю выражения до проверки возможности деления на ноль.
  \checkitem После обратной замены я решаю задачу относительно исходной переменной.
  \checkitem При решении системы нахожу пересечение всех полученных множеств.
\end{itemize}

%----------------------------------------------------------------------

\section*{\textcolor{darkaccent}{2. Границы промежутков}}
\sectionrule

От вида неравенства зависит, включается ли граничная точка в ответ.

\[
\begin{aligned}
6\le n\le 7
&\quad\Longrightarrow\quad n\in[6;7],\\
6<n<7
&\quad\Longrightarrow\quad n\in(6;7),\\
6<n\le 7
&\quad\Longrightarrow\quad n\in(6;7].
\end{aligned}
\]

Если дополнительно известно, что \(n\in\Z\), то:

\[
n\in[6;7]\cap\Z
\quad\Longrightarrow\quad
n\in\{6,7\}.
\]

\begin{itemize}
  \checkitem При знаках \(\le\) и \(\ge\) допустимая граница включается.
  \checkitem При знаках \(<\) и \(>\) граница исключается.
  \checkitem Выколотая по ОДЗ точка не включается даже при нестрогом неравенстве.
\end{itemize}

%----------------------------------------------------------------------

\section*{\textcolor{darkaccent}{3. Область допустимых значений логарифма}}
\sectionrule

Выражение
\[
\log_{a(x)} b(x)
\]
определено при одновременном выполнении условий
\[
\boxed{
a(x)>0,\qquad
a(x)\ne1,\qquad
b(x)>0.
}
\]

Для дроби
\[
\frac{f(x)}{\log_{a(x)}b(x)}
\]
добавляется условие
\[
\log_{a(x)}b(x)\ne0
\quad\Longleftrightarrow\quad
b(x)\ne1.
\]

Некоторые ограничения выполняются автоматически. Например,
\[
4+x^2>0
\]
при любом \(x\in\R\), поскольку \(x^2\ge0\). Аналогично,
\[
9\cdot3^x>0
\]
при всех действительных \(x\).

\begin{itemize}
  \checkitem Подлогарифмическое выражение всегда строго положительно.
  \checkitem Основание логарифма строго положительно и отлично от единицы.
  \checkitem Если логарифм находится в знаменателе, его значение отлично от нуля.
  \checkitem Я проверяю, не выполняется ли часть ОДЗ автоматически.
\end{itemize}

%----------------------------------------------------------------------

\section*{\textcolor{darkaccent}{4. Равенство логарифмов}}
\sectionrule

При фиксированном основании
\[
a>0,\qquad a\ne1
\]
функция \(y=\log_a x\) взаимно однозначна. Поэтому на ОДЗ
\[
\log_a f(x)=\log_a g(x)
\quad\Longleftrightarrow\quad
f(x)=g(x).
\]

\subsection*{\textcolor{accent}{Пример: логарифмическое уравнение с тригонометрией}}

Решим уравнение
\[
\log_2(\sin2x)=\log_2(\sqrt2\cos x).
\]

Сначала выпишем ОДЗ:
\[
\sin2x>0,\qquad \sqrt2\cos x>0.
\]

На ОДЗ аргументы логарифмов равны:
\[
\sin2x=\sqrt2\cos x.
\]

Поскольку
\[
\sin2x=2\sin x\cos x,
\]
получаем
\[
2\sin x\cos x-\sqrt2\cos x=0,
\]
\[
\cos x\left(2\sin x-\sqrt2\right)=0.
\]

Первая ветвь,
\[
\cos x=0,
\]
противоречит условию \(\cos x>0\). Остаётся
\[
\sin x=\frac{\sqrt2}{2}.
\]

Отсюда
\[
x=\frac{\pi}{4}+2\pi n
\quad\text{или}\quad
x=\frac{3\pi}{4}+2\pi n,
\qquad n\in\Z.
\]

Условию \(\cos x>0\) соответствует только первая серия:
\[
\boxed{x=\frac{\pi}{4}+2\pi n,\qquad n\in\Z.}
\]

\begin{itemize}
  \checkitem Я не отбрасываю логарифмы до записи ОДЗ.
  \checkitem Если аргументы логарифмов равны, одно из одинаковых по смыслу
        ограничений можно считать следствием другого.
  \checkitem Перед делением на \(\cos x\) я отдельно рассматриваю случай
        \(\cos x=0\).
  \checkitem Найденные серии корней проверяю по ОДЗ.
\end{itemize}

%----------------------------------------------------------------------

\section*{\textcolor{darkaccent}{5. Рационализация логарифмических неравенств}}
\sectionrule

Для
\[
a(x)>0,\qquad a(x)\ne1,\qquad f(x)>0
\]
справедлива равносильность
\[
\boxed{
\log_{a(x)}f(x)>0
\quad\Longleftrightarrow\quad
\bigl(a(x)-1\bigr)\bigl(f(x)-1\bigr)>0.
}
\]

Аналогично,
\[
\boxed{
\log_{a(x)}f(x)<0
\quad\Longleftrightarrow\quad
\bigl(a(x)-1\bigr)\bigl(f(x)-1\bigr)<0.
}
\]

Для сравнения двух логарифмов с одинаковым переменным основанием:
\[
\log_{a(x)}f(x)>\log_{a(x)}g(x)
\]
на ОДЗ равносильно
\[
\boxed{
\bigl(a(x)-1\bigr)\bigl(f(x)-g(x)\bigr)>0.
}
\]

Множитель \(a(x)-1\) кодирует поведение логарифмической функции:

\begin{itemize}
  \item при \(a(x)>1\) функция возрастает;
  \item при \(0<a(x)<1\) функция убывает.
\end{itemize}

Поэтому при переменном основании нельзя просто сравнить аргументы
и сохранить знак неравенства.

\subsection*{\textcolor{accent}{Обязательная формулировка перехода}}

Корректная запись в решении:

\[
\text{«На ОДЗ знак исходного выражения совпадает со знаком }
\bigl(a(x)-1\bigr)\bigl(f(x)-g(x)\bigr)\text{».}
\]

\begin{itemize}
  \checkitem Перед рационализацией я полностью выписываю ОДЗ.
  \checkitem Я учитываю множитель \(a(x)-1\).
  \checkitem Я сохраняю все выколотые точки после перехода к рациональному выражению.
  \checkitem После рационализации применяю метод интервалов.
\end{itemize}

%----------------------------------------------------------------------

\section*{\textcolor{darkaccent}{6. Логарифмические дроби}}
\sectionrule

Для дроби
\[
\frac{\log_a f(x)}{\log_a g(x)}
\]
нужно проверить:
\[
f(x)>0,\qquad
g(x)>0,\qquad
g(x)\ne1.
\]

При фиксированном основании \(a>1\) знаки логарифмов определяются так:
\[
\log_a u
\begin{cases}
<0,&0<u<1,\\
=0,&u=1,\\
>0,&u>1.
\end{cases}
\]

При \(0<a<1\) знаки меняются на противоположные.

\important{Типичная ошибка:}
\[
\log_a(uv)\ne\log_a u\cdot\log_a v.
\]

Верная формула:
\[
\boxed{\log_a(uv)=\log_a u+\log_a v}
\]
при \(u>0\) и \(v>0\).

Также:
\[
\log_a\frac{u}{v}=\log_a u-\log_a v,
\qquad
\log_a(u^k)=k\log_a u.
\]

\begin{itemize}
  \checkitem Знаменатель дроби не может равняться нулю.
  \checkitem Я не превращаю логарифм произведения в произведение логарифмов.
  \checkitem При переходе к рациональному выражению сохраняю исходное ОДЗ.
\end{itemize}

%----------------------------------------------------------------------

\section*{\textcolor{darkaccent}{7. Метод интервалов}}
\sectionrule

Для решения рационального неравенства:

\begin{enumerate}
  \item Перенесите всё в одну сторону.
  \item Разложите числитель и знаменатель на множители.
  \item Найдите нули числителя и знаменателя.
  \item Отметьте критические точки на числовой прямой.
  \item Определите знак на каждом промежутке.
  \item Выберите промежутки нужного знака.
  \item Учтите строгость неравенства и ОДЗ.
\end{enumerate}

Если корень множителя имеет нечётную кратность, знак при переходе
через него меняется. При чётной кратности знак сохраняется.

\[
\sqrt2\approx1{,}41,\qquad
\sqrt3\approx1{,}73,\qquad
\sqrt5\approx2{,}24.
\]

Эти приближения помогают быстро расположить числа на оси. Например,
\[
3-\sqrt2>0,
\]
поскольку \(\sqrt2<3\).

\begin{center}
\begin{tikzpicture}[x=1.35cm,y=1cm]
  \draw[->,color=textgray,line width=0.8pt] (-3.3,0)--(3.5,0)
    node[right] {$x$};

  \foreach \x/\lab in {-2/{-\sqrt2},0/0,2/{\sqrt2}} {
    \draw[color=textgray] (\x,0.10)--(\x,-0.10);
    \node[below=3pt] at (\x,0) {$\lab$};
  }

  \node[color=darkaccent] at (-2.7,0.45) {$+$};
  \node[color=darkaccent] at (-1,0.45) {$-$};
  \node[color=darkaccent] at (1,0.45) {$-$};
  \node[color=darkaccent] at (2.7,0.45) {$+$};
\end{tikzpicture}
\end{center}

\important{Правило системы:} если точка исключена хотя бы одним
условием системы, она не входит в итоговое пересечение.

%----------------------------------------------------------------------

\section*{\textcolor{darkaccent}{8. Показательные уравнения}}
\sectionrule

Если основания равны и удовлетворяют условиям
\[
a>0,\qquad a\ne1,
\]
то
\[
a^{f(x)}=a^{g(x)}
\quad\Longleftrightarrow\quad
f(x)=g(x).
\]

\subsection*{\textcolor{accent}{Пример}}

Решим уравнение
\[
7^{2\cos x}=49^{\sin2x}.
\]

Представим \(49\) как \(7^2\):
\[
7^{2\cos x}
=
\left(7^2\right)^{\sin2x}
=
7^{2\sin2x}.
\]

Следовательно,
\[
2\cos x=2\sin2x,
\]
\[
\cos x=\sin2x=2\sin x\cos x.
\]

Получаем
\[
\cos x(1-2\sin x)=0.
\]

Отсюда
\[
\cos x=0
\quad\text{или}\quad
\sin x=\frac12.
\]

Ответ:
\[
\boxed{
x=\frac{\pi}{2}+\pi k
\quad\text{или}\quad
x=\frac{\pi}{6}+2\pi n
\quad\text{или}\quad
x=\frac{5\pi}{6}+2\pi n,
\quad k,n\in\Z.
}
\]

Альтернативный путь состоит в логарифмировании обеих частей.
Поскольку обе части положительны, можно применить \(\log_7\):
\[
\log_7\left(7^{2\cos x}\right)
=
\log_7\left(49^{\sin2x}\right).
\]
Результат будет тем же:
\[
2\cos x=2\sin2x.
\]

\begin{itemize}
  \checkitem Сначала пытаюсь привести обе части к одному основанию.
  \checkitem При возведении степени в степень перемножаю показатели:
        \((a^m)^n=a^{mn}\).
  \checkitem Логарифмирование допустимо, когда обе части уравнения положительны.
  \checkitem Разные корректные способы должны приводить к одному множеству корней.
\end{itemize}

%----------------------------------------------------------------------

\section*{\textcolor{darkaccent}{9. Замена переменной}}
\sectionrule

Если выражение содержит \(x^4\) и \(x^2\), полезна замена
\[
t=x^2.
\]

Обязательно фиксируем ограничение:
\[
\boxed{t\ge0.}
\]

Равенство \(t=0\) допустимо, поскольку при \(x=0\) имеем \(x^2=0\).

\subsection*{\textcolor{accent}{Пример с логарифмической степенью}}

Решим неравенство
\[
2^{\log_2(4+x^2)}+x^4-10\ge0.
\]

Аргумент логарифма положителен автоматически:
\[
4+x^2>0.
\]

Используем формулу
\[
a^{\log_a u}=u
\quad
(a>0,\ a\ne1,\ u>0).
\]

Тогда
\[
4+x^2+x^4-10\ge0,
\]
\[
x^4+x^2-6\ge0.
\]

Выполним замену
\[
t=x^2,\qquad t\ge0.
\]

Получаем
\[
t^2+t-6\ge0,
\]
\[
(t+3)(t-2)\ge0.
\]

Для квадратного неравенства:
\[
t\le-3
\quad\text{или}\quad
t\ge2.
\]

Ветка \(t\le-3\) противоречит условию \(t\ge0\). Следовательно,
\[
t\ge2.
\]

Возвращаемся к переменной \(x\):
\[
x^2\ge2.
\]

Итог:
\[
\boxed{
x\in(-\infty;-\sqrt2]\cup[\sqrt2;+\infty).
}
\]

\begin{itemize}
  \checkitem При замене \(t=x^2\) я пишу \(t\ge0\).
  \checkitem Я нахожу все корни уравнения относительно \(t\).
  \checkitem Корни, противоречащие условию \(t\ge0\), исключаю с пояснением.
  \checkitem После решения относительно \(t\) выполняю обратную замену.
\end{itemize}

%----------------------------------------------------------------------

\section*{\textcolor{darkaccent}{10. Критические ошибки}}
\sectionrule

\begin{tabularx}{\linewidth}{@{}p{0.35\linewidth}X@{}}
\toprule
\textbf{\textcolor{darkaccent}{Ошибка}} &
\textbf{\textcolor{darkaccent}{Корректное действие}}\\
\midrule

Отбросить логарифмы до записи ОДЗ &
Сначала выписать все ограничения, затем выполнить равносильный переход.\\

Записать \(\log_a(uv)=\log_a u\cdot\log_a v\) &
Использовать формулу
\(\log_a(uv)=\log_a u+\log_a v\).\\

Разделить на \(\cos x\) &
Сначала отдельно рассмотреть случай \(\cos x=0\).\\

При \(t=x^2\) написать \(t>0\) &
Записать \(t\ge0\), поскольку \(x\) может равняться нулю.\\

Сравнить аргументы при переменном основании &
Учесть знак множителя \(a(x)-1\).\\

Включить нуль знаменателя в ответ &
Все нули знаменателя всегда исключаются.\\

Забыть вернуться к \(x\) после замены &
Выполнить обратную замену и решить полученное условие относительно \(x\).\\

Механически чередовать знаки &
Проверить кратность каждого корня.\\

\bottomrule
\end{tabularx}

%----------------------------------------------------------------------
% Тренировочный блок
%----------------------------------------------------------------------

\clearpage
\section*{\textcolor{darkaccent}{Тренировочный блок}}
\sectionrule

Выполните упражнения, обязательно записывая ОДЗ и обоснования
равносильных переходов.

\subsection*{\textcolor{accent}{Средний уровень}}

\begin{enumerate}
  \item Решите уравнение
  \[
  \log_3(x+2)=\log_3(2x-1).
  \]

  \item Решите показательное уравнение
  \[
  5^{x-1}=25.
  \]

  \item Решите неравенство
  \[
  x^4-5x^2+4\ge0.
  \]
\end{enumerate}

\subsection*{\textcolor{accent}{Уровень выше среднего}}

\begin{enumerate}[resume]
  \item Решите уравнение
  \[
  \log_2(\sin2x)=\log_2(\cos x).
  \]

  \item Решите неравенство
  \[
  \log_{x+2}(x^2-1)>0.
  \]

  \item Решите неравенство
  \[
  \frac{\log_2(x-1)}{\log_2(x+1)}\le0.
  \]

  \item Решите уравнение
  \[
  3^{2\cos x}=9^{\sin x}.
  \]

  \item Решите неравенство
  \[
  2^{\log_2(x^2+1)}+x^4-7\ge0.
  \]

  \item Решите неравенство
  \[
  \log_{x+1}(x^2-3x+3)\le0.
  \]
\end{enumerate}

\subsection*{\textcolor{accent}{Сложный уровень}}

\begin{enumerate}[resume]
  \item Решите неравенство
  \[
  \frac{\log_2(x^2-1)}{\log_2(x+2)}\le0.
  \]
\end{enumerate}

\clearpage
\section*{\textcolor{darkaccent}{Ответы}}
\sectionrule

\begin{longtable}{@{}p{0.08\linewidth}p{0.86\linewidth}@{}}
\toprule
\textbf{\textcolor{darkaccent}{№}} &
\textbf{\textcolor{darkaccent}{Ответ}}\\
\midrule
\endfirsthead

\toprule
\textbf{\textcolor{darkaccent}{№}} &
\textbf{\textcolor{darkaccent}{Ответ}}\\
\midrule
\endhead

1 &
\[
x=3.
\]
\\

2 &
\[
x=3.
\]
\\

3 &
\[
x\in(-\infty;-2]\cup[-1;1]\cup[2;+\infty).
\]
\\

4 &
\[
x=\frac{\pi}{6}+2\pi n,\qquad n\in\Z.
\]
\\

5 &
\[
x\in(-\sqrt2;-1)\cup(\sqrt2;+\infty).
\]
\\

6 &
\[
x\in(1;2].
\]
\\

7 &
\[
x=\frac{\pi}{4}+\pi n,\qquad n\in\Z.
\]
\\

8 &
\[
x\in(-\infty;-\sqrt2]\cup[\sqrt2;+\infty).
\]
\\

9 &
\[
x\in(-1;0)\cup[1;2].
\]
\\

10 &
\[
x\in(-2;-\sqrt2]\cup(1;\sqrt2].
\]
\\

\bottomrule
\end{longtable}

\vfill

\begin{center}
\textcolor{softgray}{%
Проверяйте ОДЗ до преобразований и после получения ответа.}
\end{center}

\end{document}